% ============================================================================
% LANGENTWURF: Las direcciones (Wegbeschreibung)
% ============================================================================
% Sequenz 6c: Mi ciudad - Las direcciones
% Fachfarbe: #c41e3a (Karminrot)
% Tier 1 (vollständig)
% ============================================================================

\documentclass[11pt,a4paper]{article}

% ============================================================================
% PACKAGES
% ============================================================================
\usepackage[utf8]{inputenc}
\usepackage[T1]{fontenc}
\usepackage[ngerman]{babel}
\usepackage{geometry}
\usepackage{fancyhdr}
\usepackage{lastpage}
\usepackage{xcolor}
\usepackage{tcolorbox}
\usepackage{enumitem}
\usepackage{tabularx}
\usepackage{longtable}
\usepackage{array}
\usepackage{booktabs}
\usepackage{hyperref}
\usepackage{ifthen}
\usepackage{amssymb}

% ============================================================================
% KONFIGURATION
% ============================================================================

\newcommand{\plantier}{1}
\newcommand{\fachid}{spanisch}
\newcommand{\fach}{Spanisch}
\newcommand{\fachfarbe}{c41e3a}

% ============================================================================
% VARIABLEN
% ============================================================================

\newcommand{\jahrgangsstufe}{7}
\newcommand{\schulart}{Gesamtschule}
\newcommand{\stundenthema}{Las direcciones -- Wegbeschreibung}
\newcommand{\reihenthema}{Mi ciudad}
\newcommand{\stundennummer}{6c}
\newcommand{\reihenstunden}{4}
\newcommand{\unterrichtsdatum}{XX.XX.XXXX}
\newcommand{\unterrichtszeit}{XX:XX -- XX:XX}
\newcommand{\raum}{XXX}

\newcommand{\lehrkraft}{N.N.}
\newcommand{\ausbildungsstand}{Lehrkraft}

\newcommand{\lerngruppengroesse}{24}
\newcommand{\maennlich}{12}
\newcommand{\weiblich}{12}
\newcommand{\divers}{0}

\newcommand{\gerniveau}{A1}
\newcommand{\lehrwerk}{Qué pasa 1}
\newcommand{\lehrwerkseite}{S. 98--103}

% ============================================================================
% LAYOUT
% ============================================================================
\geometry{left=2.5cm, right=2.5cm, top=2.5cm, bottom=2.5cm}
\definecolor{fach}{HTML}{\fachfarbe}
\definecolor{gray}{HTML}{666666}
\definecolor{lightgray}{HTML}{f5f5f5}
\definecolor{successgreen}{HTML}{2a7a4b}
\definecolor{warningorange}{HTML}{b35c00}

\tcbuselibrary{skins,breakable}

\newtcolorbox{infobox}[1][]{
    colback=lightgray,
    colframe=fach,
    fonttitle=\bfseries,
    title=#1,
    breakable
}

\newtcolorbox{scorebox}{
    colback=white,
    colframe=fach,
    fonttitle=\bfseries,
    title=Didaktik-Score,
    breakable
}

\pagestyle{fancy}
\fancyhf{}
\fancyhead[L]{\small\textcolor{gray}{\fach{} -- Jg. \jahrgangsstufe{} -- \stundenthema}}
\fancyhead[R]{\small\textcolor{gray}{Langentwurf Tier 1}}
\fancyfoot[C]{\small\thepage{} / \pageref{LastPage}}
\renewcommand{\headrulewidth}{0.4pt}
\renewcommand{\footrulewidth}{0pt}

% ============================================================================
% DOKUMENT
% ============================================================================
\begin{document}

% ============================================================================
% DECKBLATT
% ============================================================================
\begin{titlepage}
    \centering
    \vspace*{2cm}

    {\large\textcolor{gray}{\schulart}}\\[2cm]

    {\Huge\textcolor{fach}{\textbf{Langentwurf}}}\\[0.5cm]
    {\large Tier 1 -- Vollständig}\\[2cm]

    {\LARGE\textbf{\stundenthema}}\\[0.5cm]
    {\large im Rahmen der Unterrichtsreihe}\\[0.3cm]
    {\Large\textit{\reihenthema}}\\[2cm]

    \begin{tabular}{rl}
        \textbf{Fach:} & \fach{} (\gerniveau) \\[0.3cm]
        \textbf{Klasse:} & \jahrgangsstufe \\[0.3cm]
        \textbf{Lehrwerk:} & \lehrwerk, \lehrwerkseite \\[0.3cm]
        \textbf{Datum:} & \underline{\hspace{3cm}} \\[0.3cm]
        \textbf{Zeit:} & \underline{\hspace{3cm}} (Doppelstunde) \\[0.3cm]
        \textbf{Raum:} & \underline{\hspace{3cm}} \\[0.3cm]
    \end{tabular}

    \vfill

    {\large Lehrkraft: \lehrkraft}\\[0.3cm]
    {\normalsize\textcolor{gray}{\ausbildungsstand}}\\[1cm]

\end{titlepage}

% ============================================================================
% INHALTSVERZEICHNIS
% ============================================================================
\tableofcontents
\newpage

% ============================================================================
% 1. BEDINGUNGSANALYSE
% ============================================================================
\section{Bedingungsanalyse}

\subsection{Lerngruppe}
\begin{tabular}{ll}
    Kursgröße: & \lerngruppengroesse{} SuS (\maennlich{} m / \weiblich{} w / \divers{} d) \\
    Jahrgangsstufe: & \jahrgangsstufe \\
    Sprachniveau: & \gerniveau{} (GER) \\
    Fremdsprache: & 2. Fremdsprache \\
\end{tabular}

\subsection{Voraussetzungen}
Die SuS kennen aus den vorherigen Stunden:
\begin{itemize}
    \item Orte in der Stadt (la farmacia, el supermercado, la panadería, etc.) -- 6a
    \item Verben \textit{ir} und \textit{estar} mit Präpositionen -- 6b
    \item Fragen mit \textit{¿Dónde está...?}
\end{itemize}

\subsection{Differenzierung (intern)}

\textbf{WICHTIG:} Keine Sternchen auf Schülermaterial!

\begin{tabular}{|l|p{4cm}|p{4cm}|p{4cm}|}
\hline
& \textbf{[B] Basis} & \textbf{[S] Standard} & \textbf{[E] Erweiterung} \\
\hline
\textbf{Ziel} & Berufsreife & Mittlere Reife & Gymnasium \\
\hline
\textbf{Inhalt} & Einfache Richtungsangaben verstehen & Wegbeschreibung geben & Komplexe Wegbeschreibung mit mehreren Etappen \\
\hline
\textbf{Hilfen} & Bildkarten, Wortliste & Teils vorgegeben & Nur Impulse \\
\hline
\end{tabular}

\subsection{Besonderheiten}
\begin{itemize}
    \item \textbf{LRS:} 2 SuS (Nachteilsausgleich: Zeitzugabe, vergrößerte Schrift)
    \item \textbf{DaZ:} 1 SuS (Wortschatzarbeit verstärken)
    \item \textbf{Leistungsstark:} 3 SuS (Zusatzaufgaben: Rollenspiel mit Komplikationen)
\end{itemize}

% ============================================================================
% 2. SACHANALYSE
% ============================================================================
\section{Sachanalyse}

\subsection{Sprachliche Strukturen}

\textbf{Richtungsangaben (las direcciones):}
\begin{itemize}
    \item \textbf{todo recto} -- geradeaus
    \item \textbf{a la derecha} -- nach rechts
    \item \textbf{a la izquierda} -- nach links
    \item \textbf{girar} -- abbiegen
    \item \textbf{seguir} -- weitergehen
\end{itemize}

\textbf{Ortsangaben (ubicación):}
\begin{itemize}
    \item \textbf{al lado de} -- neben
    \item \textbf{enfrente de} -- gegenüber von
    \item \textbf{cerca de} -- in der Nähe von
    \item \textbf{lejos de} -- weit weg von
    \item \textbf{entre ... y ...} -- zwischen ... und ...
\end{itemize}

\textbf{Fragen:}
\begin{itemize}
    \item \textbf{¿Dónde está...?} -- Wo ist...?
    \item \textbf{¿Cómo llego a...?} -- Wie komme ich zu...?
    \item \textbf{¿Hay un/una ... por aquí?} -- Gibt es hier in der Nähe ein/eine...?
\end{itemize}

\subsection{Kontrastive Betrachtung (Deutsch $\leftrightarrow$ Spanisch)}
\begin{itemize}
    \item \textbf{Positive Transfers:} Präpositionalstruktur ähnlich (``neben dem'' = ``al lado de'')
    \item \textbf{Interferenzen:} Deutsche Präpositionen mit Dativ/Akkusativ vs. spanische feste Präpositionen
    \item \textbf{Falsche Freunde:} ``directo'' bedeutet nicht ``direkt dahin'', sondern ``direkt/gerade'' (Adjektiv)
\end{itemize}

\subsection{Kulturelle Aspekte}
\begin{itemize}
    \item Spanische Städte: Plaza Mayor als zentraler Orientierungspunkt
    \item Höfliche Anrede bei Wegbeschreibungen: \textit{Perdone/Disculpe, ¿podría decirme...?}
    \item Typische Reaktionen: \textit{¡Claro!}, \textit{Por supuesto}
\end{itemize}

% ============================================================================
% 3. DIDAKTISCHE ANALYSE
% ============================================================================
\section{Didaktische Analyse}

\subsection{Stellung in der Sequenz}
Dies ist die \textbf{3. Doppelstunde} (6c) der Sequenz ``\reihenthema'' (4 Doppelstunden gesamt).

\begin{tabular}{|c|l|l|}
\hline
\textbf{Nr.} & \textbf{Thema} & \textbf{Schwerpunkt} \\
\hline
6a & Los lugares de la ciudad & Wortschatz: Orte \\
\hline
6b & Ir y estar & Verben + Präpositionen \\
\hline
\textbf{6c} & \textbf{Las direcciones} & \textbf{Wegbeschreibung} \\
\hline
6d & Mi ciudad -- Test & Sequenztest \\
\hline
\end{tabular}

\subsection{Kompetenzziele (Rahmenplan MV)}
\begin{itemize}
    \item \textbf{Kommunikativ:}
    \begin{itemize}
        \item Hörverstehen: Wegbeschreibungen verstehen
        \item Sprechen: Wege beschreiben, nach dem Weg fragen
    \end{itemize}
    \item \textbf{Sprachlich:}
    \begin{itemize}
        \item Richtungsvokabular anwenden
        \item Präpositionen der Lage korrekt verwenden
    \end{itemize}
    \item \textbf{Interkulturell:} Höfliche Anrede, städtisches Leben in Spanien
    \item \textbf{Methodisch:} Mit Stadtplänen arbeiten
\end{itemize}

\subsection{Legitimation}
\textbf{Rahmenplan Spanisch MV (Kl. 7-10):}
\begin{itemize}
    \item Kompetenzbereich: Kommunikative Kompetenz -- Sprechen (zusammenhängendes Sprechen)
    \item Standard: ``einfache Wegbeschreibungen geben''
    \item Themenfeld: \textit{¿Dónde y cómo vivo?}
\end{itemize}

% ============================================================================
% 4. STUNDENZIELE
% ============================================================================
\section{Stundenziele}

\subsection{Grobziel}
Die SuS erweitern ihre \textbf{kommunikative Kompetenz} im Bereich \textbf{Sprechen}, indem sie Wege in einer Stadt beschreiben und nach dem Weg fragen können.

\subsection{Feinziele (operationalisiert)}
\begin{tabular}{|c|p{8cm}|c|c|}
\hline
\textbf{Nr.} & \textbf{Die SuS können...} & \textbf{AFB} & \textbf{Diff.} \\
\hline
FZ1 & die Richtungsangaben (todo recto, a la derecha, a la izquierda) korrekt aussprechen und übersetzen & I & alle \\
\hline
FZ2 & Ortsangaben (al lado de, enfrente de, cerca de, lejos de) in Sätzen anwenden & II & alle \\
\hline
FZ3 & auf die Frage \textit{¿Cómo llego a...?} eine einfache Wegbeschreibung geben & II & [S]/[E] \\
\hline
FZ4 & einen Dialog zur Wegbeschreibung mit einem Partner führen & II & alle \\
\hline
FZ5 & eine komplexe Wegbeschreibung mit mehreren Etappen formulieren & III & [E] \\
\hline
\end{tabular}

% ============================================================================
% 5. METHODISCHE ANALYSE
% ============================================================================
\section{Methodische Analyse}

\subsection{Phasierung (Doppelstunde = 2 × 45 Min.)}
\begin{enumerate}
    \item \textbf{1. Stunde:}
    \begin{itemize}
        \item Einstieg: Stadtplan-Impuls, Problemstellung ``Wie komme ich zum Bahnhof?''
        \item Erarbeitung: Richtungsvokabular mit Gesten und Bildkarten
        \item Übung: TPR (Total Physical Response) -- Anweisungen befolgen
    \end{itemize}
    \item \textbf{2. Stunde:}
    \begin{itemize}
        \item Wiederholung: Schnelles Vokabelquiz
        \item Anwendung: Rollenspiel ``Tourist fragt nach dem Weg''
        \item Sicherung: Arbeitsblatt mit Stadtplan
    \end{itemize}
\end{enumerate}

\subsection{Medien}
\begin{itemize}
    \item Tafel/Whiteboard: Stadtplan-Skizze
    \item Lehrwerk: \lehrwerk, \lehrwerkseite
    \item Arbeitsblatt: AB-06c.pdf
    \item Lernpfad: LP-06c.html (optional, digital)
    \item Bildkarten: Richtungspfeile, Orte
    \item Audio: Hörverstehensübung (Wegbeschreibung)
\end{itemize}

% ============================================================================
% 6. VERLAUFSPLANUNG
% ============================================================================
\section{Verlaufsplanung}

\textbf{1. Stunde (45 Min.)}

\begin{longtable}{|p{1cm}|p{1.5cm}|p{4cm}|p{3cm}|p{2cm}|p{2cm}|}
\hline
\textbf{Zeit} & \textbf{Phase} & \textbf{Lehrerhandeln} & \textbf{Schülerhandeln} & \textbf{SF} & \textbf{Medien} \\
\hline
\endhead

0--5 & Einstieg & L zeigt Stadtplan, fragt: ``¿Cómo llego a la estación?'' & SuS äußern Vermutungen & Plenum & Tafel/Beamer \\
\hline

5--15 & Input & L präsentiert Richtungsangaben mit Gesten: \textit{todo recto, a la derecha, a la izquierda} & SuS sprechen nach, machen Gesten mit & Plenum & Bildkarten \\
\hline

15--25 & Übung 1 & L gibt TPR-Anweisungen: ``¡Gira a la derecha!'' & SuS führen Bewegungen aus & Plenum & -- \\
\hline

25--35 & Erweiterung & L führt Ortsangaben ein: \textit{al lado de, enfrente de, cerca de, lejos de} & SuS ordnen Bilder zu & Plenum/EA & Bildkarten \\
\hline

35--40 & Sicherung & L fragt Vokabular ab & SuS antworten & Plenum & Tafel \\
\hline

40--45 & Überleitung & Zusammenfassung, Ausblick 2. Stunde & Aufräumen, Notizen & Plenum & -- \\
\hline
\end{longtable}

\vspace{1em}

\textbf{2. Stunde (45 Min.)}

\begin{longtable}{|p{1cm}|p{1.5cm}|p{4cm}|p{3cm}|p{2cm}|p{2cm}|}
\hline
\textbf{Zeit} & \textbf{Phase} & \textbf{Lehrerhandeln} & \textbf{Schülerhandeln} & \textbf{SF} & \textbf{Medien} \\
\hline
\endhead

0--5 & Wiederholung & Schnellquiz: L zeigt Pfeile, SuS nennen Richtung & SuS antworten schnell & Plenum & Bildkarten \\
\hline

5--10 & Modell & L demonstriert Dialog mit Co-Teacher/Schüler & SuS hören zu, identifizieren Strukturen & Plenum & -- \\
\hline

10--25 & Anwendung & L verteilt Rollenkarten, moderiert & SuS üben Dialoge in PA & PA & Rollenkarten \\
\hline

25--35 & Festigung & L verteilt AB, unterstützt individuell & SuS bearbeiten AB-06c & EA & AB \\
\hline

35--40 & Sicherung & L bespricht Lösungen mit Stadtplan & SuS vergleichen, korrigieren & Plenum & Tafel \\
\hline

40--45 & Abschluss & Zusammenfassung, HA erklären & SuS notieren HA & Plenum & -- \\
\hline
\end{longtable}

\textbf{Legende:} EA = Einzelarbeit, PA = Partnerarbeit, GA = Gruppenarbeit, SF = Sozialform

% ============================================================================
% 7. ERWARTETE SCHWIERIGKEITEN
% ============================================================================
\section{Erwartete Schwierigkeiten}

\begin{tabular}{|p{5cm}|p{8cm}|}
\hline
\textbf{Schwierigkeit} & \textbf{Reaktion} \\
\hline
Verwechslung derecha/izquierda & Gesten etablieren, Eselsbrücke: ``La i de izquierda = links'' \\
\hline
Aussprache von \textit{derecha} & Chorisches Sprechen, Betonung auf /ch/ \\
\hline
Präposition ``de'' vergessen bei \textit{al lado de} & Feste Wendung als Chunk einüben \\
\hline
Zeitproblem & Puffer: Rollenspiel kürzen, AB als HA \\
\hline
Heterogenität & Differenzierte Rollenkarten ([B]/[S]/[E]) \\
\hline
\end{tabular}

% ============================================================================
% 8. TAFELBILD
% ============================================================================
\section{Tafelbild}

\begin{center}
\fbox{
\begin{minipage}{0.9\textwidth}
\centering
\textbf{Las direcciones -- Die Wegbeschreibung}

\vspace{1em}

\begin{tabular}{|c|c|}
\hline
\textbf{Richtungen} & \textbf{Ortsangaben} \\
\hline
todo recto -- geradeaus & al lado de -- neben \\
a la derecha -- nach rechts & enfrente de -- gegenüber \\
a la izquierda -- nach links & cerca de -- in der Nähe \\
girar -- abbiegen & lejos de -- weit weg \\
\hline
\end{tabular}

\vspace{1em}

\textbf{Fragen:}\\
¿Dónde está la farmacia? -- Wo ist die Apotheke?\\
¿Cómo llego a la estación? -- Wie komme ich zum Bahnhof?

\vspace{1em}

[Stadtplan-Skizze]

\end{minipage}
}
\end{center}

% ============================================================================
% 9. HAUSAUFGABE
% ============================================================================
\section{Hausaufgabe}

\begin{itemize}
    \item Lehrwerk S. 102, Aufgabe 5 (Wegbeschreibung schreiben)
    \item Vokabeln lernen: Richtungsangaben und Ortsangaben
    \item Optional: Lernpfad LP-06c.html bearbeiten
\end{itemize}

% ============================================================================
% 10. ANLAGEN
% ============================================================================
\section{Anlagen}

\begin{enumerate}
    \item Arbeitsblatt (AB-06c.pdf)
    \item Musterlösung (ML-06c.pdf)
    \item Lehrerhinweise (LH-06c.pdf)
    \item Lernpfad (LP-06c.html)
    \item Rollenkarten für Dialoge
    \item Bildkarten: Richtungspfeile
    \item Stadtplan (vergrößert für Tafel)
\end{enumerate}

% ============================================================================
% 11. DIDAKTIK-SCORE
% ============================================================================
\newpage
\section{Didaktik-Score}

\begin{scorebox}
\textbf{Bewertung der didaktischen Qualität nach 10 Dimensionen}\\
\textbf{Ziel-Score: $\geq$ 9.0 (Premium-Qualität)}

\vspace{1em}
\begin{tabular}{|c|l|c|c|p{4.5cm}|}
\hline
\textbf{\#} & \textbf{Dimension} & \textbf{Gew.} & \textbf{Score} & \textbf{Notiz} \\
\hline
1 & Differenzierung & 15\% & 9/10 & Drei Niveaus, Rollenkarten differenziert \\
\hline
2 & Taxonomie (AFB) & 10\% & 9/10 & AFB I-III abgedeckt, gute Verteilung \\
\hline
3 & Lernziel-Alignment & 10\% & 10/10 & Passt zu RP MV, Themenfeld korrekt \\
\hline
4 & Aktivierung & 10\% & 10/10 & TPR, Rollenspiel = hohe Aktivierung \\
\hline
5 & Scaffolding & 10\% & 9/10 & Gesten, Bildkarten, Chunks \\
\hline
6 & Feedback-Qualität & 10\% & 9/10 & Sofortiges Feedback im Rollenspiel \\
\hline
7 & Sprachliche Klarheit & 10\% & 9/10 & Altersgerecht, klare Strukturen \\
\hline
8 & Motivation \& Relevanz & 10\% & 9/10 & Alltagsrelevanz: Urlaub, Austausch \\
\hline
9 & Inklusion (UDL) & 10\% & 9/10 & Visuell + auditiv + kinästhetisch \\
\hline
10 & Praktikabilität & 5\% & 9/10 & Zeitrahmen realistisch \\
\hline
\multicolumn{3}{|r|}{\textbf{GESAMT:}} & \textbf{9.2/10} & \\
\hline
\end{tabular}

\vspace{1em}
\textbf{Bewertungsschwellen:}
\begin{itemize}
    \item[\checkmark] \textcolor{successgreen}{$\geq$ 9.0: \textbf{FREIGABE} (Premium-Qualität)}
    \item[$\Box$] \textcolor{warningorange}{8.0--8.9: Iteration empfohlen}
    \item[$\Box$] 6.0--7.9: Überarbeitung erforderlich
    \item[$\Box$] $<$ 6.0: Grundlegende Neukonzeption
\end{itemize}
\end{scorebox}

\end{document}
